% GENERATED FILE. Edit canonical lessons, then run npm run generate:books.
% canonical-source-hash: fnv1a64:3028907a43bcabee
% canonical-lessons: GU-C302-bambo, GU-C302-bujhavvu, GU-C302-dordu, GU-C302-ogalvu, GU-C302-bhadko

\chapter{The Fire Engine}
\label{ch:gu-the-fire-engine}
\hlchaptermodality{302}

\begin{chapteropening}
\textbf{By the end of this chapter:} \emph{I can say how the fire was put out.}

Complete the last lesson of chapter 302: I can say how the fire was put out.
\end{chapteropening}

\section[\texorpdfstring{\gu{બંબો} (bambo)}{બંબો (bambo)}]{\gu{બંબો} (bambo) — a fire engine}

\label{lesson:GU-C302-bambo}



\begin{quote}
Before the new one: say the Gujarati for to burst, then the Gujarati for a bang.
\end{quote}

\subsection*{You'll want to know: \gu{બંબો}}
\textbf{\gu{બંબો}} — \emph{bambo} — \textquotedblleft{}a fire engine\textquotedblright{}.

\begin{quote}
\gu{પાંચ} \gu{મિનિટમાં} \gu{બંબો} \gu{આવી} \gu{ગયો}.
\end{quote}

\begin{quote}
\emph{pā̃c miniṭmā̃ bambo āvī gayo.}
\end{quote}

\begin{quote}
The fire engine arrived within five minutes.
\end{quote}

A masculine noun.

\begin{cousinweb}
\textbf{\gu{બંબો}} (\emph{bambo}) is probably from Portuguese \emph{bomba}, \textquotedblleft{}a pump\textquotedblright{}: the fire engine began as a pump on wheels. The firemen are \textbf{\gu{બંબાવાળા}}, the men of the pump.
\end{cousinweb}

\begin{grammarlens}[title={What you've built}]
One more word for telling what happened.
\end{grammarlens}

\subsection*{Your turn}
\begin{itemize}
\raggedright
  \item \emph{Say it:} \emph{bambo}
  \item \emph{Say it:} \emph{bambo}, once more
  \item \emph{Say it:} \emph{pā̃c miniṭmā̃ bambo āvī gayo.} — The fire engine arrived within five minutes.
  \item \emph{Recall:} say the Gujarati for to burst, then the Gujarati for a bang, then say \emph{bambo} again
\end{itemize}

\begin{tcolorbox}[breakable,colback=teal!4,colframe=teal!35!black,title={Before you move on}]
What does \gu{બંબો} mean? (\textquotedblleft{}A fire engine\textquotedblright{}.) Say it once more.
\end{tcolorbox}
\section[\texorpdfstring{\gu{બુઝાવવું} (bujhāvvũ)}{બુઝાવવું (bujhāvvũ)}]{\gu{બુઝાવવું} (bujhāvvũ) — to extinguish}

\label{lesson:GU-C302-bujhavvu}



\begin{quote}
Before the new one: say the Gujarati for a bang, then the Gujarati for a fire engine.
\end{quote}

\subsection*{You'll want to know: \gu{બુઝાવવું}}
\textbf{\gu{બુઝાવવું}} — \emph{bujhāvvũ} — \textquotedblleft{}to extinguish\textquotedblright{}.

\begin{quote}
\gu{લોકોએ} \gu{પાણી} \gu{નાખીને} \gu{આગ} \gu{બુઝાવી}.
\end{quote}

\begin{quote}
\emph{lokoe pāṇī nākhīne āg bujhāvī.}
\end{quote}

\begin{quote}
People poured water on the fire and put it out.
\end{quote}

Past, with \gu{તેણે} (he or she …): \gu{તેણે} \gu{આગ} \gu{બુઝાવી}, \emph{teṇe āg bujhāvī}.

\begin{cousinweb}
\textbf{\gu{બુઝાવવું}} (\emph{bujhāvvũ}) belongs with \emph{bujhāvũ}, to go out, which CDIAL lists under \emph{*vijjhāyati}, \textquotedblleft{}is burnt out, is extinguished\textquotedblright{}, a form it reconstructs; the doing verb continues that entry's causative line, \emph{*vijjhāpayati} (CDIAL 11703): to make a fire go out.
\end{cousinweb}

\begin{grammarlens}[title={What you've built}]
One more word for telling what happened.
\end{grammarlens}

\subsection*{Your turn}
\begin{itemize}
\raggedright
  \item \emph{Say it:} \emph{bujhāvvũ}
  \item \emph{Say it:} \emph{bujhāvvũ}, once more
  \item \emph{Say it:} \emph{lokoe pāṇī nākhīne āg bujhāvī.} — People poured water on the fire and put it out.
  \item \emph{Recall:} say the Gujarati for a bang, then the Gujarati for a fire engine, then say \emph{bujhāvvũ} again
\end{itemize}

\begin{tcolorbox}[breakable,colback=teal!4,colframe=teal!35!black,title={Before you move on}]
What does \gu{બુઝાવવું} mean? (\textquotedblleft{}To extinguish\textquotedblright{}.) Say it once more.
\end{tcolorbox}
\section[\texorpdfstring{\gu{દોરડું} (dorḍũ)}{દોરડું (dorḍũ)}]{\gu{દોરડું} (dorḍũ) — a rope}

\label{lesson:GU-C302-dordu}



\begin{quote}
Before the new one: say the Gujarati for a fire engine, then the Gujarati for to extinguish.
\end{quote}

\subsection*{You'll want to know: \gu{દોરડું}}
\textbf{\gu{દોરડું}} — \emph{dorḍũ} — \textquotedblleft{}a rope\textquotedblright{}.

\begin{quote}
\gu{તેમણે} \gu{દોરડાથી} \gu{એક} \gu{બાળકને} \gu{નીચે} \gu{ઉતાર્યું}.
\end{quote}

\begin{quote}
\emph{temṇe dorḍāthī ek bāḷakne nīce utāryũ.}
\end{quote}

\begin{quote}
They let a child down on a rope.
\end{quote}

A neuter noun.

\begin{cousinweb}
\textbf{\gu{દોરડું}} (\emph{dorḍũ}) belongs with \textbf{\gu{દોરો}}, a thread: both go back to Sanskrit \emph{davara}, \textquotedblleft{}string\textquotedblright{}, where CDIAL lists Gujarati \emph{doro}, \textquotedblleft{}string, rope\textquotedblright{} (CDIAL 6225). The ending \emph{-ḍũ} makes the thicker thing, a rope.
\end{cousinweb}

\begin{grammarlens}[title={What you've built}]
One more word for telling what happened.
\end{grammarlens}

\subsection*{Your turn}
\begin{itemize}
\raggedright
  \item \emph{Say it:} \emph{dorḍũ}
  \item \emph{Say it:} \emph{dorḍũ}, once more
  \item \emph{Say it:} \emph{temṇe dorḍāthī ek bāḷakne nīce utāryũ.} — They let a child down on a rope.
  \item \emph{Recall:} say the Gujarati for a fire engine, then the Gujarati for to extinguish, then say \emph{dorḍũ} again
\end{itemize}

\begin{tcolorbox}[breakable,colback=teal!4,colframe=teal!35!black,title={Before you move on}]
What does \gu{દોરડું} mean? (\textquotedblleft{}A rope\textquotedblright{}.) Say it once more.
\end{tcolorbox}
\section[\texorpdfstring{\gu{ઓગળવું} (ogaḷvũ)}{ઓગળવું (ogaḷvũ)}]{\gu{ઓગળવું} (ogaḷvũ) — to melt}

\label{lesson:GU-C302-ogalvu}



\begin{quote}
Before the new one: say the Gujarati for to extinguish, then the Gujarati for a rope.
\end{quote}

\subsection*{You'll want to know: \gu{ઓગળવું}}
\textbf{\gu{ઓગળવું}} — \emph{ogaḷvũ} — \textquotedblleft{}to melt\textquotedblright{}.

\begin{quote}
\gu{તડકામાં} \gu{માખણ} \gu{ઓગળી} \gu{ગયું}.
\end{quote}

\begin{quote}
\emph{taḍkāmā̃ mākhaṇ ogḷī gayũ.}
\end{quote}

\begin{quote}
The butter melted in the sun.
\end{quote}

Past (it …): \gu{તે} \gu{બરફ} \gu{ઓગળી} \gu{ગયો}, \emph{te baraph ogḷī gayo}.

\begin{cousinweb}
\textbf{\gu{ઓગળવું}} (\emph{ogaḷvũ}) is inherited from Sanskrit \emph{avagalati}, \textquotedblleft{}trickles down\textquotedblright{}, and CDIAL lists the Gujarati verb, \textquotedblleft{}to trickle down, melt\textquotedblright{} (CDIAL 742). Ice and butter \textbf{\gu{ઓગળે} \gu{છે}}.
\end{cousinweb}

\begin{grammarlens}[title={What you've built}]
One more word for telling what happened.
\end{grammarlens}

\subsection*{Your turn}
\begin{itemize}
\raggedright
  \item \emph{Say it:} \emph{ogaḷvũ}
  \item \emph{Say it:} \emph{ogaḷvũ}, once more
  \item \emph{Say it:} \emph{taḍkāmā̃ mākhaṇ ogḷī gayũ.} — The butter melted in the sun.
  \item \emph{Recall:} say the Gujarati for to extinguish, then the Gujarati for a rope, then say \emph{ogaḷvũ} again
\end{itemize}

\begin{tcolorbox}[breakable,colback=teal!4,colframe=teal!35!black,title={Before you move on}]
What does \gu{ઓગળવું} mean? (\textquotedblleft{}To melt\textquotedblright{}.) Say it once more.
\end{tcolorbox}
\section[\texorpdfstring{\gu{ભડકો} (bhaḍko)}{ભડકો (bhaḍko)}]{\gu{ભડકો} (bhaḍko) — a blaze; a sudden flare of flame}

\label{lesson:GU-C302-bhadko}



\begin{quote}
Before the new one: say the Gujarati for a rope, then the Gujarati for to melt.
\end{quote}

\subsection*{You'll want to know: \gu{ભડકો}}
\textbf{\gu{ભડકો}} — \emph{bhaḍko} — \textquotedblleft{}a blaze; a sudden flare of flame\textquotedblright{}.

\begin{quote}
\gu{તપેલીમાં} \gu{તેલ} \gu{ખૂબ} \gu{ગરમ} \gu{થયું} \gu{અને} \gu{મોટો} \gu{ભડકો} \gu{થયો}.
\end{quote}

\begin{quote}
\emph{tapelīmā̃ tel khūb garam thayũ ane moṭo bhaḍko thayo.}
\end{quote}

\begin{quote}
The oil in the pan got too hot and flared up in a big blaze.
\end{quote}

A masculine noun.

\begin{cousinweb}
comes from \emph{*bhaṭ}, \textquotedblleft{}a sudden movement or noise\textquotedblright{}, a form CDIAL reconstructs, under which it lists Gujarati \emph{bhaṛkũ}, \textquotedblleft{}a blaze\textquotedblright{} (CDIAL 9365). A \textbf{\gu{ભડકો}} is flame that leaps up all at once.
\end{cousinweb}

\begin{grammarlens}[title={What you've built}]
One more word for telling what happened.
\end{grammarlens}

\subsection*{Your turn}
\begin{itemize}
\raggedright
  \item \emph{Say it:} \emph{bhaḍko}
  \item \emph{Say it:} \emph{bhaḍko}, once more
  \item \emph{Say it:} \emph{tapelīmā̃ tel khūb garam thayũ ane moṭo bhaḍko thayo.} — The oil in the pan got too hot and flared up in a big blaze.
  \item \emph{Recall:} say the Gujarati for a rope, then the Gujarati for to melt, then say \emph{bhaḍko} again
\end{itemize}

\begin{tcolorbox}[breakable,colback=teal!4,colframe=teal!35!black,title={Before you move on}]
What does \gu{ભડકો} mean? (\textquotedblleft{}A blaze; a sudden flare of flame\textquotedblright{}.) Say it once more.
\end{tcolorbox}
